\begin{corollary}\label{fga-cor-8-2-6-2}\dcref{8.2.6.2}\source{fga-explained:page-0200}\uses{fga-cor-8-2-4}
Under the hypotheses of Corollary~\ref{fga-cor-8-2-4}, put
$B=\bigoplus_{n\ge0}f_*\mathcal O_X(n)$ (with the grading induced by the
chosen ample ideal).  For every $q$, the graded $B$-module
\[
\bigoplus_{n\ge0}R^qf_*(\mathcal F\otimes\mathcal I^n)
\]
is finitely generated, and there is $n_0$ such that multiplication by
$\mathcal I$ induces isomorphisms
$R^qf_*(\mathcal F\otimes\mathcal I^n)\simeq
R^qf_*(\mathcal F\otimes\mathcal I^{n+1})$ for all $n\ge n_0$.
\end{corollary}
\begin{proof}
Use the exact sequence for successive powers of $\mathcal I$ and the comparison
theorem to obtain a graded finite $B$-module controlling the kernels and images.
Finite generation over the noetherian ring $B$ gives stabilization of the graded
images; the exact sequence then gives the asserted stabilization of the direct
images and the finite generation of the displayed module.
\end{proof}

\begin{lemma}\label{fga-lem-8-2-7-4}\dcref{8.2.7.4}\source{fga-explained:page-0202}\uses{}
For an exact sequence of inverse systems of $A$-modules
\[
0\to L'_\bullet\to L_\bullet\to L''_\bullet\to0,
\]
if $L''_\bullet$ satisfies the Artin--Rees Mittag--Leffler condition (ARML),
then so does $L'_\bullet$; if both $L'_\bullet$ and $L''_\bullet$ satisfy
ARML, then so does $L_\bullet$.
\end{lemma}
\begin{proof}
The transition images in the left system are kernels of the corresponding images
in the middle system.  Eventual stabilization in the quotient therefore implies
eventual stabilization in the submodule.  Applying this observation twice to the
short exact sequence proves the two assertions.
\end{proof}

\begin{corollary}\label{fga-cor-8-2-13}\dcref{8.2.13}\source{fga-explained:page-0204}\uses{fga-thm-8-2-12}
Under the hypotheses of Theorem~\ref{fga-thm-8-2-12}, for every $y\in Y$ the
connected components of $X_y$ are in bijection with the points of
$Y'=\operatorname{Spec}_Y(f_*\mathcal O_X)$ over $y$, equivalently with the
maximal ideals of the finite local $\mathcal O_{Y,y}$-algebra
$(f_*\mathcal O_X)_y$.
\end{corollary}
\begin{proof}
The Stein factorization gives $X\to Y'\to Y$ with finite second map and
geometrically connected fibres for the first.  The fibre $Y'_y$ is finite
and discrete, and the connected components of $X_y$ are exactly its inverse
images.  This gives the stated bijection.
\end{proof}

\begin{corollary}\label{fga-cor-8-2-14}\dcref{8.2.14}\source{fga-explained:page-0204}\uses{fga-thm-8-2-12}
Let $f:X\to Y$ be proper and surjective between integral noetherian schemes,
with $Y$ normal.  If the generic fibre is geometrically connected, then every
fibre of $f$ is geometrically connected.
\end{corollary}
\begin{proof}
In the Stein factorization $X\to Y'\to Y$, the generic fibre hypothesis says
that the extension of function fields defining $Y'$ is radicial.  Normality of
$Y$ implies that the normalization in this radicial extension is local at every
point, so each fibre of $Y'\to Y$ has one point after every field extension.
Corollary~\ref{fga-cor-8-2-13} then gives geometric connectedness of every
fibre of $f$.
\end{proof}

\begin{corollary}\label{fga-cor-8-2-15}\dcref{8.2.15}\source{fga-explained:page-0204}\uses{fga-thm-8-2-12}
Under the hypotheses of Theorem~\ref{fga-thm-8-2-12}, a point $x\in X$ is open
and closed in its fibre if and only if the fibre of $X\to Y'$ over $f'(x)$ is
the singleton $\{x\}$.  The locus $U$ of such points is open, $f'(U)$ is open
in $Y'$, and $U\to f'(U)$ is an isomorphism.
\end{corollary}
\begin{proof}
The finite morphism $Y'\to Y$ has discrete fibres, so isolation can be checked
after replacing $Y'$ by the point $f'(x)$.  Properness makes the image of the
complement of an affine neighbourhood of $x$ closed; shrinking the target gives
an affine neighbourhood whose inverse image is that neighbourhood.  The map on
coordinate rings is an isomorphism by the singleton-fibre condition and
$f_*\mathcal O_X=\mathcal O_{Y'}$.
\end{proof}

\begin{corollary}\label{fga-cor-8-2-16}\dcref{8.2.16}\source{fga-explained:page-0204}\uses{}
A proper quasi-finite morphism of locally noetherian schemes is finite.
\end{corollary}
\begin{proof}
Quasi-finiteness gives finite fibres and hence, by the preceding isolated-point
criterion and Zariski's Main Theorem, an affine neighbourhood of every target
point on which the morphism is finite.  Finiteness is affine-local on the target,
so the morphism is finite globally.
\end{proof}

\begin{remark}\label{fga-rem-8-3-5}\dcref{8.3.5}\source{fga-explained:page-0208}\uses{fga-thm-8-3-2}
The comparison result extends to quasi-compact, quasi-separated morphisms by
replacing the finite Cech complex for an affine cover with a hyper-Cech complex.
The tor-independence hypothesis in the basic statement cannot in general be
removed by this argument.
\end{remark}

\begin{remark}\label{fga-rem-8-3-11-1}\dcref{8.3.11.1}\source{fga-explained:page-0212}\uses{}
In the cohomological-flatness criterion, condition (a) for $q=1$ is automatic;
thus $Rf_*\mathcal F$ has perfect amplitude in $[0,1]$.  Condition (b) for
$q=0$ is likewise automatic, so $\mathcal F$ is cohomologically flat in degree
zero.
\end{remark}

\begin{remark}\label{fga-rem-8-3-11-2}\dcref{8.3.11.2}\source{fga-explained:page-0212}\uses{fga-thm-8-3-8}
If $H^1(X_y,\mathcal F/\mathfrak m_y\mathcal F)=0$, formal functions imply
$R^1f_*\mathcal F_y=0$ and the base-change map for $f_*\mathcal F$ is
surjective.  In a neighbourhood of $y$, $f_*\mathcal F$ is then locally free
of finite type and commutes with arbitrary base change.  More generally, if
$Rf_*\mathcal F$ has perfect amplitude in $[q+1,b]$, the dual perfect complex
$K=R\mathcal Hom(R^qf_*\mathcal F,\mathcal O_Y)$ gives a quasi-coherent sheaf
$Q=\mathcal H^{-q}(K)$ and natural exchange isomorphisms
\[
R^{b-q}f_*(\mathcal M\otimes\mathcal F)\simeq
\mathcal Hom_Y(Q,\mathcal M),
\]
compatible with base change.
\end{remark}

\begin{remark}\label{fga-rem-8-4-3-2}\dcref{8.4.3.2}\source{fga-explained:page-0215}\uses{fga-lem-8-4-3-1}
The algebraization lemma remains valid for an adic formal scheme $\mathfrak X$
whose special fibre is proper and equipped with an ample invertible sheaf.  A
uniform twist makes the restriction maps on sections surjective at all levels,
so a coherent sheaf with proper support algebraizes; the quasi-projective and
general cases follow by extension by zero and noetherian induction.
\end{remark}

\begin{remark}\label{fga-rem-8-4-8}\dcref{8.4.8}\source{fga-explained:page-0217}\uses{fga-cor-8-4-7}
Properness is essential for algebraization.  For $X=\operatorname{Spec}\mathbb Z$
with its $0$-adic formal completion, continuous homomorphisms of formal power
series rings are restricted power series, whereas algebraic morphisms are
polynomials; the natural map from algebraic to formal morphisms is not
surjective.
\end{remark}

\begin{corollary}\label{fga-cor-8-5-6}\dcref{8.5.6}\source{fga-explained:page-0222}\uses{fga-thm-8-4-2,fga-cor-8-4-7}
Let $\mathfrak X$ be a proper flat adic locally noetherian formal scheme over
$S$.  If it is the completion of a proper scheme, then that scheme is flat over
$S$.  If $H^1(X_0,\mathcal O_{X_0})=0$, every line bundle on the special fibre
$X_0$ lifts uniquely up to isomorphism to $\mathfrak X$.  If $X_0$ is projective
and an ample line bundle lifts, then $\mathfrak X$ algebraizes to a proper flat
projective scheme with an ample line bundle inducing the given formal data.
\end{corollary}
\begin{proof}
Flatness follows from the local flatness criterion applied to the thickenings.
The obstruction to lifting a line bundle from level $n$ to $n+1$ lies in
$H^2(X_0,\mathcal O_{X_0}\otimes I_n)$ and the ambiguity in $H^1$; the stated
vanishing gives existence and uniqueness.  Algebraization of the line bundle and
formal scheme is Corollary~\ref{fga-cor-8-4-7}.
\end{proof}

\begin{theorem}\label{fga-the-8-5-9}\dcref{8.5.9}\source{fga-explained:page-0223}\uses{fga-def-6-4-2}
Let $Y$ be smooth over $S$, let $j:X_0\hookrightarrow X$ be a square-zero
thickening with ideal $J$, and let $g:X_0\to Y$ be an $S$-morphism.  There is
an obstruction
\[
o(g,j)\in H^1(X_0,J\otimes g^*T_{Y/S})
\]
whose vanishing is equivalent to the existence of an extension $X\to Y$; when
it vanishes, extensions form a torsor under
$H^0(X_0,J\otimes g^*T_{Y/S})$.

If $S_0\hookrightarrow S$ has square-zero ideal $I$ and $X_0$ is smooth over
$S_0$, there is an obstruction
$o(X_0,S)\in H^2(X_0,f_0^*I\otimes T_{X_0/S_0})$.  Its vanishing is equivalent
to existence of a deformation over $S$; isomorphism classes form a torsor under
$H^1(X_0,f_0^*I\otimes T_{X_0/S_0})$, and automorphisms of a fixed deformation
are $H^0$ of the same sheaf.  In particular, etale $X_0/S_0$ has a unique
deformation.
\end{theorem}
\begin{proof}
Smoothness gives local liftings.  Two local liftings differ by an
$S$-derivation, hence by a section of $J\otimes g^*T_{Y/S}$; their Cech
differences define the first obstruction class and its torsor of extensions.
For deformations of $X_0$, local smooth liftings form a gerbe banded by
$f_0^*I\otimes T_{X_0/S_0}$; its class is the displayed $H^2$ obstruction,
and neutralizations and automorphisms give the $H^1$ and $H^0$ torsors.
\end{proof}

\begin{remark}\label{fga-rem-8-5-13}\dcref{8.5.13}\source{fga-explained:page-0227}\uses{fga-thm-8-5-12}
By Artin approximation, the specialization and lifting conclusions remain true
when the complete base ring is replaced by a henselian local noetherian ring.
These results are the input for proper base change in etale cohomology.
\end{remark}

\begin{theorem}\label{fga-the-8-5-31}\dcref{8.5.31}\source{fga-explained:page-0236}\uses{fga-def-6-4-2}
Let $X_0\hookrightarrow X$ be a square-zero thickening with ideal $J$ and let
$g:X_0\to Y$ be an $S$-morphism.  The obstruction to extending $g$ to $X$ is
\[
o(g,j)\in\operatorname{Ext}^1_{X_0}(Lg^*L_{Y/S},J).
\]
It vanishes exactly when an extension exists; when it vanishes, extensions form
a torsor under $\operatorname{Ext}^0_{X_0}(Lg^*L_{Y/S},J)$.  For a square-zero
base thickening $S_0\hookrightarrow S$, the obstruction to deforming a flat
$X_0/S_0$ lies in $\operatorname{Ext}^2(L_{X_0/S_0},f_0^*I)$, with deformation
and automorphism groups given by $\operatorname{Ext}^1$ and $\operatorname{Ext}^0$.
\end{theorem}
\begin{proof}
The transitivity triangle for the cotangent complex identifies liftings with
splittings of the corresponding square-zero extension.  Applying
$R\mathcal Hom(-,J)$ gives the obstruction boundary in degree one and the
torsor of splittings in degree zero.  Applying the same argument to the
structural morphism $X_0\to S_0$ shifts the obstruction to degree two and gives
the stated deformation and automorphism groups.
\end{proof}

\begin{remark}\label{fga-rem-8-5-33}\dcref{8.5.33}\source{fga-explained:page-0238}\uses{fga-cor-8-5-32}
Under the hypotheses of Corollary~\ref{fga-cor-8-5-32}, completion identifies
the cotangent complex of the algebraized scheme with the completed cotangent
complex of the formal scheme.
\end{remark}

\begin{remark}\label{fga-rem-8-6-8}\dcref{8.6.8}\source{fga-explained:page-0241}\uses{}
In Serre's characteristic-$p$ construction, let $A$ be a complete local
noetherian ring with residue field $k$ and let $\mathfrak X$ be a formal versal
deformation of the special fibre.  The fixed-point estimate
$r+\dim(Q_0)<n$ forces the image of $p$ in $A$ to be nilpotent.  Equivalently,
the formal deformation has characteristic $p^N$ for some $N$; Serre's argument
shows in the relevant range that one may take $N=1$.
\end{remark}
